\section{Experimental Setup}
\label{sec:setup}

This chapter outlines the complete experimental testbed configuration
required to evaluate Xen and KVM under mixed-criticality workloads. To
ensure reliable, predictable, and reproducible latency measurements, the
host environment must be strictly tailored for real-time execution.

The following sections detail the step-by-step preparation of the
system, starting with the installation and tuning of the host operating
system, the compilation of a fully preemptible Linux kernel
(\texttt{PREEMPT\_RT}), and the deployment of the respective
hypervisors. We then detail both hypervisors' configuration, focusing on
the differences between the two.

\subsection{Kernel configuration and real-time
tuning}\label{kernel-configuration-and-real-time-tuning}

To establish a baseline for our performance comparison, we also required
the standard, non-real-time Linux kernel version 6.18.35.

\begin{itemize}
\item
  First, we downloaded the official kernel source code from the
  \href{https://cdn.kernel.org/pub/linux/kernel/}{Linux Kernel Archive}.
\item
  We extracted the archive using the following command:

\begin{Shaded}
\begin{Highlighting}[]
\FunctionTok{tar}\NormalTok{ {-}xzvf \textasciitilde{}/Downloads/linux{-}6.18.35.tar.gz {-}C \textasciitilde{}/}
\end{Highlighting}
\end{Shaded}
\item
  Next, we installed the necessary dependencies to build the kernel:

\begin{Shaded}
\begin{Highlighting}[]
\FunctionTok{sudo}\NormalTok{ apt {-}y install libncurses{-}dev gawk flex bison openssl libssl{-}dev dkms libelf{-}dev libudev{-}dev libpci{-}dev libiberty{-}dev autoconf llvm qtcreator qtbase5{-}dev qt5{-}qmake cmake}
\end{Highlighting}
\end{Shaded}
\item
  We navigated into the Linux build tree and copied the configuration
  file from the currently running system:

\begin{Shaded}
\begin{Highlighting}[]
\BuiltInTok{cd}\NormalTok{ \textasciitilde{}/linux{-}6.18.35}
\FunctionTok{cp}\NormalTok{ /boot/config{-}}\VariableTok{$(}\FunctionTok{uname}\NormalTok{ {-}r}\VariableTok{)}\NormalTok{ .config}
\end{Highlighting}
\end{Shaded}
\item
  To fix potential configuration issues arising from the kernel version
  mismatch, we updated the configuration:

\begin{Shaded}
\begin{Highlighting}[]
\FunctionTok{make}\NormalTok{ olddefconfig}
\end{Highlighting}
\end{Shaded}
\item
  To streamline the build by compiling only the currently loaded
  modules, we ran:

\begin{Shaded}
\begin{Highlighting}[]
\FunctionTok{make}\NormalTok{ localmodconfig}
\end{Highlighting}
\end{Shaded}
\end{itemize}

\subsubsection{Patching Linux with
PREEMPT\_RT}\label{patching-linux-with-preempt_rt}

To achieve true real-time determinism after installation, we must move
beyond the default \texttt{PREEMPT\_DYNAMIC} schema. While
\texttt{PREEMPT\_DYNAMIC} allows the kernel to dynamically determine
preemption modes (e.g., \emph{none}, \emph{voluntary}, or \emph{full}),
it is not designed for real-time workloads and lacks hard guarantees for
interrupt latency and thread scheduling.

For the Real-Time kernel, we needed to follow some additional steps:

\begin{itemize}
\item
  We downloaded the matching PREEMPT\_RT patch from the
  \href{https://wiki.linuxfoundation.org/realtime/start}{Linux
  Foundation Real-Time Wiki} and extracted it:

\begin{Shaded}
\begin{Highlighting}[]
\FunctionTok{gunzip}\NormalTok{ {-}c \textasciitilde{}/Downloads/patch{-}6.18.35{-}rt5.patch.gz }\OperatorTok{\textgreater{}}\NormalTok{ \textasciitilde{}/patch{-}6.18.35{-}rt5.patch}
\end{Highlighting}
\end{Shaded}
\item
  Returning to the build tree, we applied the patch and opened the
  configuration GUI:

\begin{Shaded}
\begin{Highlighting}[]
\BuiltInTok{cd}\NormalTok{ \textasciitilde{}/linux{-}6.18.35}
\FunctionTok{patch}\NormalTok{ {-}p1 }\OperatorTok{\textless{}}\NormalTok{ ../patch{-}6.18.35{-}rt5.patch}
\FunctionTok{make}\NormalTok{ xconfig}
\end{Highlighting}
\end{Shaded}
\item
  In the configuration menu, we navigated to General Setup
  -\textgreater{} Preemption Model and set it to Fully Preemptible
  Kernel (RT).
\item
  We opened the configuration menu again (\texttt{make\ xconfig}) to
  manually disable specific features to reduce latency, strictly in the
  following order:

  \begin{itemize}
  \item
    \texttt{CONFIG\_SCHED\_MC\_PRIO} (\textbf{Processor type and
    features} -\textgreater{} \textbf{Multi-core scheduler support})
  \item
    \texttt{CONFIG\_CPU\_FREQ} (\textbf{Power management and ACPI
    options} -\textgreater{} \textbf{CPU Frequency scaling})
  \item
    \texttt{CONFIG\_STACKPROTECTOR}
  \item
    \texttt{CONFIG\_APM}
  \item
    \texttt{CONFIG\_ACPI\_PROCESSOR} (\textbf{Power management and ACPI
    options} -\textgreater{} \textbf{ACPI (Advanced Configuration and
    Power Interface)})
  \item
    \texttt{CONFIG\_CPU\_IDLE} (\textbf{Power management and ACPI
    options} -\textgreater{} \textbf{CPU idle PM support})
  \item
    \textbf{Simultaneous Multi-threading} (if supported by the hardware)
  \end{itemize}
\item
  We also ensured NVMe support was enabled:

  \begin{itemize}
  \item
    \texttt{CONFIG\_BLK\_DEV\_NVME} (\textbf{Device Drivers}
    -\textgreater{} \textbf{NVME Support} -\textgreater{} \textbf{NVM
    Express block device})
  \end{itemize}
\item
  Under \textbf{Processor type and features}, we fine-tuned the settings
  for our specific hardware architecture and saved the configuration.
\end{itemize}

\subsubsection{Compiling the kernel}\label{compiling-the-kernel}

\begin{itemize}
\item
  For Ubuntu specifically, it is necessary to clear the Canonical
  certificates (\texttt{canonical.pem}) to avoid build failures. From
  the Linux build tree, we executed:

\begin{Shaded}
\begin{Highlighting}[]
\FunctionTok{sudo}\NormalTok{ scripts/config {-}{-}disable SYSTEM\_TRUSTED\_KEYS}
\FunctionTok{sudo}\NormalTok{ scripts/config {-}{-}disable SYSTEM\_REVOCATION\_KEYS}
\FunctionTok{make}\NormalTok{ olddefconfig}
\end{Highlighting}
\end{Shaded}
\item
  Finally, we built and installed the kernel and its modules:

\begin{Shaded}
\begin{Highlighting}[]
\FunctionTok{sudo}\NormalTok{ make {-}j25}
\FunctionTok{sudo}\NormalTok{ make modules\_install}
\FunctionTok{sudo}\NormalTok{ make install}
\end{Highlighting}
\end{Shaded}
\end{itemize}

\subsubsection{Hardware Tuning
  (BIOS/UEFI)}\label{hardware-tuning-biosuefi}

The configurations in this section are highly hardware-specific and will
vary depending on the motherboard manufacturer and CPU vendor. To ensure
predictable performance and minimize latency spikes for real-time
workloads, it is crucial to disable dynamic frequency scaling and deep
power-saving states directly at the firmware level.

Below are the specific configuration steps applied to our testbed, which
features an MSI motherboard and an AMD processor:

\begin{itemize}
\item
  \textbf{Disable Precision Boost Overdrive (PBO):} Navigate to
  \texttt{Overclocking} -\textgreater{}
  \texttt{Advanced\ CPU\ Configuration} -\textgreater{}
  \texttt{AMD\ Overclocking} and set \textbf{Precision Boost Overdrive}
  to \textbf{Disabled}.
\item
  \textbf{Disable CPU Boost and Sleep States:} Navigate to
  \texttt{Overclocking} -\textgreater{}
  \texttt{Advanced\ CPU\ Configuration} -\textgreater{}
  \texttt{AMD\ CBS} and set both \textbf{Core Performance Boost} and
  \textbf{Global C-state Control} to \textbf{Disabled}.
\item
  \textbf{Set a Fixed CPU Frequency:} In the main \texttt{Overclocking}
  menu, change the \textbf{CPU Ratio} from \texttt{Auto} to a fixed
  value of \textbf{44.00}. This forces the processor to run at a
  constant clock speed, preventing the latency overhead associated with
  dynamic frequency transitions during the experiments.
\end{itemize}

\subsubsection{OS-level isolation}\label{os-level-isolation}

Given the presence of some counterintuitive results, we decided to
configure the operating system to adopt additional isolation mechanisms,
aiming to identify the possible cause of such behaviors by minimizing
the variables introduced by the OS scheduler.

\paragraph{Isolation and Configuration
Steps}\label{isolation-and-configuration-steps}

The isolation was implemented at the operating system level through
various techniques, including:

\begin{itemize}
\item
  \textbf{Kernel scheduling isolation (\texttt{isolcpus}):} Use the
  \texttt{isolcpus=} boot parameter to prevent the scheduler from
  assigning tasks to a specific set of CPUs, thereby avoiding general
  SMP balancing.
\item
  \textbf{Kernel house-keeping noise (\texttt{nohz\_full}):} Address
  kernel house-keeping noise by appending \texttt{nohz\_full=} to enable
  Tickless Mode, which reduces OS overhead on those selected CPUs.
\item
  \textbf{RCU callback offloading (\texttt{rcu\_nocbs}):} To further
  minimize latencies imposed by memory allocators in \texttt{softirq}
  contexts, apply RCU callback offloading to dedicated kernel threads
  using the \texttt{rcu\_nocbs=} parameter.
\item
  \textbf{IRQ affinity:} Configure IRQ affinity to ensure hardware
  interrupts are kept away from our isolated cores as much as possible.
\end{itemize}

The isolation was applied to both the host and the virtualized
environment. The host system was configured to avoid scheduling tasks on
the physical CPUs dedicated to running the virtual machines.
Concurrently, the guest system was also configured so that the virtual
machine's scheduler could not assign tasks to the cores reserved for
real-time applications.

The following is an example of OS-level isolation we adopted in one of
the KVM experiments. On the host, we fully isolated \textbf{CPU22} and
\textbf{CPU23} on our 24-core system, configuring GRUB by adding a
custom boot entry in \texttt{/etc/grub.d/40\_custom}:

\begin{Verbatim}[breaklines=true,breakanywhere=true]
menuentry 'Ubuntu 24.04 (6.18.35-rt5-full)'{
        echo 'Loading Linux 6.18.35-rt5 with NO_HZ, ISOLCPUS, RCU_NOCBS and no IRQ_AFFINITY on [22, 23]'
        linux   /boot/vmlinuz-6.18.35-rt5 root=UUID=8e001ea3-a450-434d-bfab-ee2e6f61c6a2 ro  nohz_full=22-23 isolcpus=22-23 rcu_nocbs=22-23 irqaffinity=0-21 quiet splash $vt_handoff
        echo    'Loading initial ramdisk ...'
        initrd  /boot/initrd.img-6.18.35-rt5
}
\end{Verbatim}

Similarly, on the guest we isolated \textbf{CPU1} with the following
configuration:

\begin{Verbatim}[breaklines=true,breakanywhere=true]
menuentry 'Ubuntu 24.04 (6.18.35-rt5-full)'{
        echo 'Loading Linux 6.18.35-rt5 with NO_HZ, ISOLCPUS, RCU_NOCBS and no IRQ_AFFINITY on 1'
        linux   /boot/vmlinuz-6.18.35-rt5 root=UUID=1fe78ac9-040a-4fe5-b3a8-8715d38d692e ro  nohz_full=1 isolcpus=1 rcu_nocbs=1 irqaffinity=0 quiet splash $vt_handoff
        echo    'Loading initial ramdisk ...'
        initrd  /boot/initrd.img-6.18.35-rt5
}
\end{Verbatim}

\subsection{KVM}\label{kvm}

Being effectively treated as a Type-2 hypervisor, KVM sits on top of an
already booted operating system. The host OS manages it similarly to a
user-space application, where the virtual CPUs (vCPUs) are scheduled as
standard host processes. Consequently, the installation process is as
straightforward as running the following command:

\begin{Shaded}
\begin{Highlighting}[]
\FunctionTok{sudo}\NormalTok{ apt {-}y install bridge{-}utils cpu{-}checker libvirt{-}clients libvirt{-}daemon qemu{-}system qemu{-}kvm virt{-}manager numad}
\end{Highlighting}
\end{Shaded}

We installed QEMU emulator version 8.2.2 (Debian 1:8.2.2+ds-0ubuntu1.17)

In order to apply the OS-level isolation techniques on the pinned
configuration, we modified the \texttt{/etc/grub.d/40\_custom} file
adding the entries in the \href{KVM/grub_cfg/40_custom_host}{GRUB host
configuration}.

Once the installation was complete, we provisioned the following guest
Virtual Machines with the same hardware specifications:

\begin{table*}
\centering
\caption{KVM Virtual Machines Configuration}
\label{tab:setup:kvm-vms}
\begin{tabularx}{\linewidth}{@{}YYY@{}}
\toprule
Name & ubuntu24.04 & ubuntu24.04-pinned\tabularnewline
\midrule
\textbf{vCPUs:} & 2 & 2\tabularnewline
\textbf{RAM:} & 4 GB & 4 GB\tabularnewline
\textbf{Isolation techniques} & Memory locking, CPU pass-through & CPU pinning, memory locking, CPU pass-through, OS-level isolation\tabularnewline
\textbf{Configuration file} & \href{KVM/vms/ubuntu24.04.xml}{ubuntu24.04.xml} & \href{KVM/vms/ubuntu24.04-pinned.xml}{ubuntu24.04-pinned.xml}\tabularnewline
\bottomrule
\end{tabularx}
\end{table*}

We installed both the 6.18.35 and the 6.18.35-rt5 kernels in the same
manner as on the host, with the exception that we did not enable the
NVMe block device support, since we will be using VirtIO. For simplicity
reasons, the two VMs share the same virtual drive.

In order to apply the OS-level isolation techniques on the pinned
configuration, we modified the \texttt{/etc/grub.d/40\_custom} file
adding the entries in the \href{KVM/grub_cfg/40_custom_guest}{GRUB guest
configuration}.

We ensured the VM vCPUs always had the required resources available and
could not be preempted by other tasks by setting the \texttt{vcpusched}
mode of both vCPUs to \texttt{SCHED\_FIFO} with priority 98.

\begin{Shaded}
\begin{Highlighting}[]
\KeywordTok{\textless{}vcpusched}\OtherTok{ vcpus=}\StringTok{\textquotesingle{}0\textquotesingle{}}\OtherTok{ scheduler=}\StringTok{\textquotesingle{}fifo\textquotesingle{}}\OtherTok{ priority=}\StringTok{\textquotesingle{}98\textquotesingle{}}\KeywordTok{/\textgreater{}}
\KeywordTok{\textless{}vcpusched}\OtherTok{ vcpus=}\StringTok{\textquotesingle{}1\textquotesingle{}}\OtherTok{ scheduler=}\StringTok{\textquotesingle{}fifo\textquotesingle{}}\OtherTok{ priority=}\StringTok{\textquotesingle{}98\textquotesingle{}}\KeywordTok{/\textgreater{}}
\end{Highlighting}
\end{Shaded}

We enabled memory locking to prevent swapping:

\begin{Shaded}
\begin{Highlighting}[]
\KeywordTok{\textless{}memoryBacking\textgreater{}}
  \KeywordTok{\textless{}locked/\textgreater{}}
\KeywordTok{\textless{}/memoryBacking\textgreater{}}
\end{Highlighting}
\end{Shaded}

We ensured CPU pass-through and correct topology mapping:

\begin{Shaded}
\begin{Highlighting}[]
\KeywordTok{\textless{}cpu}\OtherTok{ mode=}\StringTok{\textquotesingle{}host{-}passthrough\textquotesingle{}}\OtherTok{ check=}\StringTok{\textquotesingle{}none\textquotesingle{}}\KeywordTok{\textgreater{}}
  \KeywordTok{\textless{}topology}\OtherTok{ sockets=}\StringTok{\textquotesingle{}1\textquotesingle{}}\OtherTok{ dies=}\StringTok{\textquotesingle{}1\textquotesingle{}}\OtherTok{ cores=}\StringTok{\textquotesingle{}2\textquotesingle{}}\OtherTok{ threads=}\StringTok{\textquotesingle{}1\textquotesingle{}}\KeywordTok{/\textgreater{}}
\KeywordTok{\textless{}/cpu\textgreater{}}
\end{Highlighting}
\end{Shaded}

In the pinned configuration, we applied CPU pinning to make the vCPU
threads run exactly on the isolated cores:

\begin{Shaded}
\begin{Highlighting}[]
\KeywordTok{\textless{}vcpu}\OtherTok{ placement=}\StringTok{\textquotesingle{}static\textquotesingle{}}\KeywordTok{\textgreater{}}\NormalTok{2}\KeywordTok{\textless{}/vcpu\textgreater{}}
\KeywordTok{\textless{}cputune\textgreater{}}
  \KeywordTok{\textless{}vcpupin}\OtherTok{ vcpu=}\StringTok{\textquotesingle{}0\textquotesingle{}}\OtherTok{ cpuset=}\StringTok{\textquotesingle{}22\textquotesingle{}}\KeywordTok{/\textgreater{}}
  \KeywordTok{\textless{}vcpupin}\OtherTok{ vcpu=}\StringTok{\textquotesingle{}1\textquotesingle{}}\OtherTok{ cpuset=}\StringTok{\textquotesingle{}23\textquotesingle{}}\KeywordTok{/\textgreater{}} 
  \KeywordTok{\textless{}emulatorpin}\OtherTok{ cpuset=}\StringTok{\textquotesingle{}0{-}21\textquotesingle{}}\KeywordTok{/\textgreater{}}
  \KeywordTok{\textless{}iothreadpin}\OtherTok{ iothread=}\StringTok{\textquotesingle{}1\textquotesingle{}}\OtherTok{ cpuset=}\StringTok{\textquotesingle{}0{-}21\textquotesingle{}}\KeywordTok{/\textgreater{}}
\KeywordTok{\textless{}/cputune\textgreater{}}
\end{Highlighting}
\end{Shaded}

\subsubsection{TACLe Benchmark}\label{tacle-benchmark}

We provisioned the following guest Virtual Machines to run the TACLe
benchmarks:

\begin{table*}
\centering
\caption{TACLe Benchmark VMs Configuration (KVM)}
\label{tab:setup:kvm-tacle}
\begin{tabularx}{\linewidth}{@{}YYY@{}}
\toprule
Name & ubuntu24.04-tacle & ubuntu24.04-tacle-pinned\tabularnewline
\midrule
\textbf{vCPUs:} & 2 & 2\tabularnewline
\textbf{RAM:} & 4 GB & 4 GB\tabularnewline
\textbf{Isolation techniques} & Memory locking, CPU pass-through & CPU pinning, memory locking, CPU pass-through, OS-level isolation\tabularnewline
\textbf{Configuration file} & \href{KVM/vms/ubuntu24.04-tacle.xml}{ubuntu24.04-tacle.xml} & \href{KVM/vms/ubuntu24.04-tacle-pinned.xml}{ubuntu24.04-tacle-pinned.xml}\tabularnewline
\bottomrule
\end{tabularx}
\end{table*}

We configured them the same way we did for the \texttt{ubuntu24.04} and
\texttt{ubuntu24.04-pinned} VMs respectively, with the exception of the
pinning layout of the emulation and I/O threads. These four also share
the same virtual drive, since only one of them will be turned on at a
time.

For the noisy guests, we provisioned similar VMs, but with different
specs and no real-time vCPU priority assigned:

\begin{table*}
\centering
\caption{Noisy Guest VMs Configuration (KVM)}
\label{tab:setup:kvm-noisy}
\begin{tabularx}{\linewidth}{@{}YYYY@{}}
\toprule
Name & ubuntu24.04-noise-small & ubuntu24.04-noise-big & ubuntu24.04-noise-big-pinned\tabularnewline
\midrule
\textbf{vCPUs:} & 2 & 16 & 16\tabularnewline
\textbf{RAM:} & 4 GB & 16 GB & 16 GB\tabularnewline
\textbf{Isolation techniques} & Memory locking, CPU pass-through & Memory locking, CPU pass-through & CPU pinning, memory locking, CPU pass-through\tabularnewline
\textbf{Configuration file} & \href{KVM/vms/ubuntu24.04-noise-small.xml}{ubuntu24.04-noise-small.xml} & \href{KVM/vms/ubuntu24.04-noise-big.xml}{ubuntu24.04-noise-big.xml} & \href{KVM/vms/ubuntu24.04-noise-big-pinned.xml}{ubuntu24.04-noise-big-pinned.xml}\tabularnewline
\bottomrule
\end{tabularx}
\end{table*}

These three also share a different virtual drive, that does not require
any modification but the installation of \texttt{stress-ng}:

\begin{Shaded}
\begin{Highlighting}[]
\FunctionTok{sudo}\NormalTok{ apt {-}y install stress{-}ng}
\end{Highlighting}
\end{Shaded}

The host configuration for the TACLe experiment differs from the
cyclictest one in a subtle but important way. In the cyclictest
experiments, the host was booted with \texttt{isolcpus=22-23} to fully
remove those two cores from the kernel's load balancer, ensuring that
QEMU's vCPU threads would be the only workload on those cores. This
approach works precisely because there is only one VM and its placement
is static.

For the TACLe experiment, however, we need the Linux scheduler to remain
active across CPUs 4--23, because both the TACLe VM and the noisy VM
must be free to float across that entire pool. Using
\texttt{isolcpus=4-23} would remove those cores from the load balancer
entirely, preventing KVM from distributing the two VMs' vCPU threads
across them --- exactly the opposite of what we need.

Instead, we boot the host using a dedicated GRUB entry that applies
\texttt{nohz\_full=4-23}, \texttt{rcu\_nocbs=4-23}, and
\texttt{irqaffinity=0-3} to reduce kernel noise on the VM cores, but
deliberately omits \texttt{isolcpus}:

\begin{Verbatim}[breaklines=true,breakanywhere=true]
menuentry 'Ubuntu 24.04 (6.18.35-rt5 TACLe)'{
        echo 'Loading Linux 6.18.35-rt5 with NO_HZ, RCU_NOCBS and no IRQ_AFFINITY on [4, 23]'
        linux   /boot/vmlinuz-6.18.35-rt5 root=UUID=8e001ea3-a450-434d-bfab-ee2e6f61c6a2 ro  nohz_full=4-23 rcu_nocbs=4-23 irqaffinity=0-3 quiet splash $vt_handoff
        initrd  /boot/initrd.img-6.18.35-rt5
}
\end{Verbatim}

To prevent host user-space services from wandering onto CPUs 4--23
without disabling the load balancer, we confine the host's systemd
cgroup to CPUs 0--3 by adding the following line to
\texttt{/etc/systemd/system.conf}:

\begin{Shaded}
\begin{Highlighting}[]
\DataTypeTok{CPUAffinity}\OtherTok{=}\DecValTok{0{-}3}
\end{Highlighting}
\end{Shaded}

This setting is applied at boot (PID 1 reads \texttt{system.conf} on
startup) and removed at the end of the experiment via
\texttt{sudo\ systemctl\ daemon-reexec}, which re-executes systemd
in-place without restarting any services, immediately restoring normal
scheduling.

\subsection{Xen}\label{xen}

We installed the Xen hypervisor 4.17.3 via the
\texttt{xen-hypervisor-amd64} package. This process automatically
generated the necessary GRUB bootloader entries to boot the Ubuntu
system as \textbf{Dom0} (the privileged management domain).

\begin{Shaded}
\begin{Highlighting}[]
\FunctionTok{sudo}\NormalTok{ apt {-}y install xen{-}hypervisor{-}amd64}
\end{Highlighting}
\end{Shaded}

During our initial boot tests, we encountered severe instability with
the Graphical User Interface (GUI). Specifically, the \texttt{nouveau}
open-source drivers---often relied upon for NVIDIA GPU
compatibility---failed to initialize correctly on our testbed. Further
investigation suggested that graphical drivers generally exhibit poor
stability when running under Xen Dom0, a behavior observed across
different hardware configurations. To bypass this limitation, we opted
for a headless setup and managed the host via SSH.

To streamline the provisioning of subsequent DomUs, we installed the
\texttt{xen-tools} package:

\begin{Shaded}
\begin{Highlighting}[]
\FunctionTok{sudo}\NormalTok{ apt {-}y install xen{-}tools}
\end{Highlighting}
\end{Shaded}

Since our guests required dedicated block storage, we resized the
existing LVM (Logical Volume Manager) partition hosting the Ubuntu
installation to carve out a new logical volume exclusively dedicated to
the VMs. We performed this using the following steps:

\begin{Shaded}
\begin{Highlighting}[]
\FunctionTok{sudo}\NormalTok{ lvreduce {-}{-}resizefs {-}{-}size {-}65G /dev/ubuntu{-}vg/root}
\FunctionTok{sudo}\NormalTok{ lvcreate {-}L 65G {-}n ubuntu{-}24.04{-}domU ubuntu{-}vg}
\end{Highlighting}
\end{Shaded}

To speed up the setup of the DomU, we decided to clone the Dom0
partition and assign it a new UUID:

\begin{Shaded}
\begin{Highlighting}[]
\FunctionTok{sudo}\NormalTok{ dd if=/dev/nvme0n1p4 of=/dev/ubuntu{-}vg/ubuntu{-}24.04{-}domU status=progress}
\FunctionTok{sudo}\NormalTok{ e2fsck {-}f /dev/ubuntu{-}vg/ubuntu{-}24.04{-}domU}
\FunctionTok{sudo}\NormalTok{ tune2fs {-}U random /dev/ubuntu{-}vg/ubuntu{-}24.04{-}domU}
\end{Highlighting}
\end{Shaded}

We designed different configuration files to provision our DomU
instances, with the same hardware specifications. The differences
between the configurations are the underlying kernel used to boot them
and how vCPU allocation is performed.

\begin{table*}
\centering
\caption{Xen DomU Configurations}
\label{tab:setup:xen-domu}
\begin{tabularx}{\linewidth}{@{}YYYYY@{}}
\toprule
Name & ubuntu-24.04-linux-6.18.35-nrt-hvm & ubuntu-24.04-linux-6.18.35-nrt-hvm-pinned & ubuntu-24.04-linux-6.18.35-rt5-hvm & ubuntu-24.04-linux-6.18.35-rt5-hvm-pinned\tabularnewline
\midrule
\textbf{vCPUs:} & 2 & 2 & 2 & 2\tabularnewline
\textbf{RAM:} & 4 GB & 4 GB & 4 GB & 4 GB\tabularnewline
\textbf{Isolation techniques} & None & CPU pinning & None & CPU pinning\tabularnewline
\textbf{Configuration file} & \href{Xen/domains/ubuntu-24.04-linux-6.18.35-nrt-hvm.conf}{ubuntu-24.04-linux-6.18.35-nrt-hvm.conf} & \href{Xen/domains/ubuntu-24.04-linux-6.18.35-nrt-hvm-pinned.conf}{ubuntu-24.04-linux-6.18.35-nrt-hvm-pinned.conf} & \href{Xen/domains/ubuntu-24.04-linux-6.18.35-rt5-hvm.conf}{ubuntu-24.04-linux-6.18.35-rt5-hvm.conf} & \href{Xen/domains/ubuntu-24.04-linux-6.18.35-rt5-hvm-pinned.conf}{ubuntu-24.04-linux-6.18.35-rt5-hvm-pinned.conf}\tabularnewline
\bottomrule
\end{tabularx}
\end{table*}

These different configurations share the same LVM partition, as only one
of them will be turned on at a time.

\subsubsection{Problems with
PREEMPT\_RT}\label{problems-with-preempt_rt}

During the initial setup phase, we attempted to boot Xen using the same
real-time kernel---compiled with the previously described
instructions---as the Dom0 kernel. We tested various kernel versions
across different Linux distributions and experimented with several
combinations of the tuning parameters mentioned earlier (specifically:
\texttt{CONFIG\_SCHED\_MC\_PRIO}, \texttt{CONFIG\_CPU\_FREQ},
\texttt{CONFIG\_STACKPROTECTOR}, \texttt{CONFIG\_APM},
\texttt{CONFIG\_ACPI\_PROCESSOR}, and \texttt{CONFIG\_CPU\_IDLE}).

Furthermore, we applied a wide array of Xen and kernel command-line boot
parameters that are traditionally recommended for resolving boot hangs
and hardware initialization issues. However, none of these mitigations
proved successful.

We also ensured that all the necessary configuration flags required to
run the kernel as a Xen Dom0 were strictly enabled. We tried:

\begin{itemize}
\item
  \texttt{swiotlb=35536,force}
\item
  \texttt{iommu=pt}
\item
  \texttt{console=hvc0}
\item
  \texttt{console=tty0}
\item
  \texttt{earlyprintk=xen}
\item
  \texttt{apci=off\ noapic\ pci=nomsi}
\item
  \texttt{noirqbalance}
\item
  \texttt{nomodeset}
\end{itemize}

Despite these extensive troubleshooting efforts, we observed that
enabling the ``Fully Preemptible Kernel (RT)'' option consistently
caused severe boot incompatibilities with the Xen hypervisor. The boot
sequence systematically stalled even before the initialization of the
logging daemons (such as \texttt{systemd-journald}).

Coupled with the graphical driver issues discussed previously, debugging
this behavior proved to be a formidable challenge. The system most
likely dropped into an \texttt{initramfs} recovery shell, which remained
completely inaccessible to us in our headless setup.

Consequently, we decided to leave the ``Fully Preemptible Kernel''
option disabled for the Dom0 kernel. Instead, we opted for the
``Low-Latency'' scheduling model, which guaranteed a reliable boot
process while still offering better responsiveness compared to the
standard generic kernel.

\subsubsection{Resource partitioning and
NULL-scheduler}\label{resource-partitioning-and-null-scheduler}

In order to compare the effects of the scheduler choice on Xen, we
swapped the default Credit2 scheduler with a pinned configuration,
effectively using an offline scheduler. This is done to assess the
current effects of the issues identified by the previous analyses of
Abeni and Faggioli, such as the priority inversion via QEMU and the
\texttt{TIMER\_SLOP} limitation. In order to do so, we changed the Dom0
GRUB boot configuration at \texttt{/etc/grub.d/40\_custom} adding the
entries shown in \href{Xen/grub_cfg/40_custom}{this GRUB configuration
file}.

This configuration restricts Xen to using only the first 22 pCPUs for
Dom0. Also, we configured the pinned guests to use only the pCPUs 22 and
23, so that they would have two dedicated cores with no interference
from the Dom0.

To be absolutely certain of the vCPU to pCPU fixed mapping, we also
executed the following commands after the VM booted:

\begin{Shaded}
\begin{Highlighting}[]
\FunctionTok{sudo}\NormalTok{ xl vcpu{-}pin ubuntu{-}24.04{-}linux{-}6.18.35{-}rt5 0 22}
\FunctionTok{sudo}\NormalTok{ xl vcpu{-}pin ubuntu{-}24.04{-}linux{-}6.18.35{-}rt5 1 23}
\end{Highlighting}
\end{Shaded}

\subsubsection{TACLe Benchmark}\label{tacle-benchmark-1}

For the TACLe benchmarks, we used the Low-Latency Dom0 and Real-Time
DomU setup, and pinned Dom0 to only use the first four pCPUs, leaving
the others to both the noisy guest and the benchmark guest. This is to
separate the cores dedicated to the privileged domain from the other
guests, ensuring that only stress generated by the benchmarks and the
noisy VM is reflected in the results.

We designed the following configurations to run the TACLe benchmarks:

\begin{table*}
\centering
\caption{TACLe Benchmark Configurations (Xen)}
\label{tab:setup:xen-tacle}
\begin{tabularx}{\linewidth}{@{}YYY@{}}
\toprule
Name & ubuntu-24.04-linux-6.18.35-rt5-hvm-tacle & ubuntu-24.04-linux-6.18.35-rt5-hvm-tacle-pinned\tabularnewline
\midrule
\textbf{vCPUs:} & 2 & 2\tabularnewline
\textbf{RAM:} & 4 GB & 4 GB\tabularnewline
\textbf{Isolation techniques} & CPU scheduling range (4-23) & CPU pinning\tabularnewline
\textbf{Configuration file} & \href{Xen/domains/ubuntu-24.04-linux-6.18.35-rt5-hvm-tacle.conf}{ubuntu-24.04-linux-6.18.35-rt5-hvm-tacle.conf} & \href{Xen/domains/ubuntu-24.04-linux-6.18.35-rt5-hvm-tacle-pinned.conf}{ubuntu-24.04-linux-6.18.35-rt5-hvm-tacle-pinned.conf}\tabularnewline
\bottomrule
\end{tabularx}
\end{table*}

The only difference between these configurations and the previous ones
is the range the vCPUs will be scheduled in. These too share the same
LVM partition assigned to the other DomUs.

For the noisy guests, we defined similar configurations, but with
different specs:

\begin{table*}
\centering
\caption{Noisy Guest Configurations (Xen)}
\label{tab:setup:xen-noisy}
\begin{tabularx}{\linewidth}{@{}YYYY@{}}
\toprule
Name & ubuntu-24.04-linux-6.18.35-nrt-hvm-noisyguest-small & ubuntu-24.04-linux-6.18.35-nrt-hvm-noisyguest-big & ubuntu-24.04-linux-6.18.35-nrt-hvm-noisyguest-big-pinned\tabularnewline
\midrule
\textbf{vCPUs:} & 2 & 16 & 16\tabularnewline
\textbf{RAM:} & 4 GB & 16 GB & 16 GB\tabularnewline
\textbf{Isolation techniques} & CPU scheduling range (4-23) & CPU scheduling range (4-23) & CPU pinning\tabularnewline
\textbf{Configuration file} & \href{Xen/domains/ubuntu-24.04-linux-6.18.35-nrt-hvm-noisyguest-small.conf}{ubuntu-24.04-linux-6.18.35-nrt-hvm-noisyguest-small.conf} & \href{Xen/domains/ubuntu-24.04-linux-6.18.35-nrt-hvm-noisyguest-big.conf}{ubuntu-24.04-linux-6.18.35-nrt-hvm-noisyguest-big.conf} & \href{Xen/domains/ubuntu-24.04-linux-6.18.35-nrt-hvm-noisyguest-big-pinned.conf}{ubuntu-24.04-linux-6.18.35-nrt-hvm-noisyguest-big-pinned.conf}\tabularnewline
\bottomrule
\end{tabularx}
\end{table*}

These three also share a different LVM partition, that does not require
any modification but the installation of \texttt{stress-ng}:

\begin{Shaded}
\begin{Highlighting}[]
\FunctionTok{sudo}\NormalTok{ apt {-}y install stress{-}ng}
\end{Highlighting}
\end{Shaded}

\paragraph{TACLe and stress on Dom0}\label{tacle-and-stress-on-dom0}

There is also a variation of the above configuration, where the stress
is generated on Dom0 itself. In this case, there is no need to limit the
scheduling range of the DomU. The configuration file is
\href{Xen/domains/ubuntu-24.04-linux-6.18.35-rt5-hvm-tacle-dom0noise.conf}{this
one}. The pinned configuration, by definition, does not change.

\subsubsection{PV and PVH DomUs}\label{pv-and-pvh-domus}

To test different virtualization technologies, we created some PV and
PVH configuration variants of the previous real-time DomUs. As we will
discuss in the Xen analysis section, the PV DomUs don't support
\texttt{PREEMPT\_RT} patched kernels with the Fully Preemptible Kernel
option enabled (in a similar fashion to Dom0, it being a PV domain too),
thus we used the Low-Latency variant of the kernel for these
configurations:

\begin{table*}
\centering
\caption{PV and PVH DomU Configurations}
\label{tab:setup:xen-pvh}
\begin{tabularx}{\linewidth}{@{}YYYYY@{}}
\toprule
Name & ubuntu-24.04-linux-6.18.35-rt5-ll-pv & ubuntu-24.04-linux-6.18.35-rt5-ll-pv-pinned & ubuntu-24.04-linux-6.18.35-rt5-pvh & ubuntu-24.04-linux-6.18.35-rt5-pvh-pinned\tabularnewline
\midrule
\textbf{vCPUs:} & 2 & 2 & 2 & 2\tabularnewline
\textbf{RAM:} & 4 GB & 4 GB & 4 GB & 4 GB\tabularnewline
\textbf{Isolation techniques} & None & CPU pinning & None & CPU pinning\tabularnewline
\textbf{Configuration file} & \href{Xen/domains/ubuntu-24.04-linux-6.18.35-rt5-ll-pv.conf}{ubuntu-24.04-linux-6.18.35-rt5-ll-pv.conf} & \href{Xen/domains/ubuntu-24.04-linux-6.18.35-rt5-ll-pv-pinned.conf}{ubuntu-24.04-linux-6.18.35-rt5-ll-pv-pinned.conf} & \href{Xen/domains/ubuntu-24.04-linux-6.18.35-rt5-pvh.conf}{ubuntu-24.04-linux-6.18.35-rt5-pvh.conf} & \href{Xen/domains/ubuntu-24.04-linux-6.18.35-rt5-pvh-pinned.conf}{ubuntu-24.04-linux-6.18.35-rt5-pvh-pinned.conf}\tabularnewline
\bottomrule
\end{tabularx}
\end{table*}

All these configurations use the same LVM partition as the other HVM
configurations, and do not require any other specific setup.

\subsection{Benchmark automation}\label{benchmark-automation}

In order to streamline the testing procedure, a series of supplementary
configuration steps were required. Since these procedures apply
identically across both virtualization platforms, the following
terminology will be used for simplicity: the term \textbf{Host} will
refer collectively to the KVM Host and the Xen Dom0, while the term
\textbf{Guest} will denote both the KVM virtual machine and the Xen
DomU.

First, we installed \texttt{openssh-server} and \texttt{expect} on the
Host to orchestrate the benchmarks from a remote machine and interact
with the Guest OS via the \texttt{xl} and \texttt{virsh} consoles.

\begin{Shaded}
\begin{Highlighting}[]
\FunctionTok{sudo}\NormalTok{ apt {-}y install openssh{-}server expect}
\end{Highlighting}
\end{Shaded}

To enable passwordless SSH access, we generated an SSH key pair on the
orchestrator machine and appended the public key to the
\texttt{authorized\_keys} file on the Host:

\begin{Shaded}
\begin{Highlighting}[]
\FunctionTok{ssh{-}keygen}\NormalTok{ {-}t ed25519 {-}C }\StringTok{"xen{-}vs{-}kvm{-}orchestrator"}
\ExtensionTok{ssh{-}copy{-}id}\NormalTok{ matt@192.168.1.166}
\end{Highlighting}
\end{Shaded}

To facilitate operations requiring root privileges without manual
intervention, we added the following rule to the \texttt{/etc/sudoers}
file on the Host and Guests using \texttt{visudo}:

\begin{Verbatim}[breaklines=true,breakanywhere=true]
matt ALL=(ALL) NOPASSWD: ALL
\end{Verbatim}

To automatically boot on the appropriate kernel with the correct
settings, we modified the GRUB configuration in
\texttt{/etc/default/grub} on the Host and KVM Guests:

\begin{Verbatim}[breaklines=true,breakanywhere=true]
GRUB_DEFAULT=saved
GRUB_TIMEOUT_STYLE=menu
GRUB_TIMEOUT=5
\end{Verbatim}

For the KVM Guests, we enabled the serial console by enabling the
corresponding \texttt{systemctl} service:

\begin{Shaded}
\begin{Highlighting}[]
\FunctionTok{sudo}\NormalTok{ systemctl enable {-}{-}now serial{-}getty@ttyS0.service}
\end{Highlighting}
\end{Shaded}
